\documentclass[10pt,a4paper]{amsart}
\usepackage{amsmath,amssymb,amsthm,graphicx}
\usepackage{verbatim}
\newenvironment{reference}{\comment}{\endcomment}
%\newenvironment{reference}{}{}
\newenvironment{slogan}{\comment}{\endcomment}
\newenvironment{history}{\comment}{\endcomment}

% For commutative diagrams we use Xy-pic
\usepackage[all]{xy}

% We use 2cell for 2-commutative diagrams.
\xyoption{2cell}
% Theorem environments.
%
\theoremstyle{plain}
\newtheorem{theorem}[subsection]{Theorem}
\newtheorem{proposition}[subsection]{Proposition}
\newtheorem{lemma}[subsection]{Lemma}

\theoremstyle{definition}
\newtheorem{definition}[subsection]{Definition}
\newtheorem{example}[subsection]{Example}
\newtheorem{exercise}[subsection]{Exercise}
\newtheorem{situation}[subsection]{Situation}

\theoremstyle{remark}
\newtheorem{remark}[subsection]{Remark}
\newtheorem{remarks}[subsection]{Remarks}

\numberwithin{equation}{subsection}

% Macros
%
\def\lim{\mathop{\mathrm{lim}}\nolimits}
\def\colim{\mathop{\mathrm{colim}}\nolimits}
\def\Spec{\mathop{\mathrm{Spec}}}
\def\Hom{\mathop{\mathrm{Hom}}\nolimits}
\def\Ext{\mathop{\mathrm{Ext}}\nolimits}
\def\SheafHom{\mathop{\mathcal{H}\!\mathit{om}}\nolimits}
\def\SheafExt{\mathop{\mathcal{E}\!\mathit{xt}}\nolimits}
\def\Sch{\mathit{Sch}}
\def\Mor{\mathop{\mathrm{Mor}}\nolimits}
\def\Ob{\mathop{\mathrm{Ob}}\nolimits}
\def\Sh{\mathop{\mathit{Sh}}\nolimits}
\def\NL{\mathop{N\!L}\nolimits}
\def\CH{\mathop{\mathrm{CH}}\nolimits}
\def\proetale{{pro\text{-}\acute{e}tale}}
\def\etale{{\acute{e}tale}}
\def\QCoh{\mathit{QCoh}}
\def\Ker{\mathop{\mathrm{Ker}}}
\def\Im{\mathop{\mathrm{Im}}}
\def\Coker{\mathop{\mathrm{Coker}}}
\def\Coim{\mathop{\mathrm{Coim}}}

% Boxtimes
%
\DeclareMathSymbol{\boxtimes}{\mathbin}{AMSa}{"02}

%
% Macros for moduli stacks/spaces
%
\def\QCohstack{\mathcal{QC}\!\mathit{oh}}
\def\Cohstack{\mathcal{C}\!\mathit{oh}}
\def\Spacesstack{\mathcal{S}\!\mathit{paces}}
\def\Quotfunctor{\mathrm{Quot}}
\def\Hilbfunctor{\mathrm{Hilb}}
\def\Curvesstack{\mathcal{C}\!\mathit{urves}}
\def\Polarizedstack{\mathcal{P}\!\mathit{olarized}}
\def\Complexesstack{\mathcal{C}\!\mathit{omplexes}}
% \Pic is the operator that assigns to X its picard group, usage \Pic(X)
% \Picardstack_{X/B} denotes the Picard stack of X over B
% \Picardfunctor_{X/B} denotes the Picard functor of X over B
\def\Pic{\mathop{\mathrm{Pic}}\nolimits}
\def\Picardstack{\mathcal{P}\!\mathit{ic}}
\def\Picardfunctor{\mathrm{Pic}}
\def\Deformationcategory{\mathcal{D}\!\mathit{ef}}
\usepackage{cleveref}
\crefname{originallemma}{original lemma}{original lemmas}
\Crefname{originallemma}{Original lemma}{Original lemmas}
\makeatletter
\newlabel{lemma-essential-image-j-shriek}{{00Y1}{source}{}{sourcepointer-1}{}}
\newlabel{lemma-essential-image-j-shriek@cref}{{[originallemma][0][]00Y1}{[1][0][]source}{}{}{}}
\makeatother
\crefname{originalsection}{original section}{original sections}
\Crefname{originalsection}{Original section}{Original sections}
\makeatletter
\newlabel{section-functoriality-PSh}{{00VC}{source}{}{sourcepointer-2}{}}
\newlabel{section-functoriality-PSh@cref}{{[originalsection][0][]00VC}{[1][0][]source}{}{}{}}
\makeatother
\crefname{originalsection}{original section}{original sections}
\Crefname{originalsection}{Original section}{Original sections}
\makeatletter
\newlabel{section-continuous-functors}{{00WU}{source}{}{sourcepointer-3}{}}
\newlabel{section-continuous-functors@cref}{{[originalsection][0][]00WU}{[1][0][]source}{}{}{}}
\makeatother
\crefname{originallemma}{original lemma}{original lemmas}
\Crefname{originallemma}{Original lemma}{Original lemmas}
\makeatletter
\newlabel{lemma-technical-up}{{00WY}{source}{}{sourcepointer-4}{}}
\newlabel{lemma-technical-up@cref}{{[originallemma][0][]00WY}{[1][0][]source}{}{}{}}
\makeatother
\crefname{originallemma}{original lemma}{original lemmas}
\Crefname{originallemma}{Original lemma}{Original lemmas}
\makeatletter
\newlabel{lemma-pullback-representable-sheaf}{{04D3}{source}{}{sourcepointer-5}{}}
\newlabel{lemma-pullback-representable-sheaf@cref}{{[originallemma][0][]04D3}{[1][0][]source}{}{}{}}
\makeatother
\crefname{originallemma}{original lemma}{original lemmas}
\Crefname{originallemma}{Original lemma}{Original lemmas}
\makeatletter
\newlabel{categories-lemma-adjoint-exact}{{0038}{source}{}{sourcepointer-6}{}}
\newlabel{categories-lemma-adjoint-exact@cref}{{[originallemma][0][]0038}{[1][0][]source}{}{}{}}
\makeatother
\crefname{originallemma}{original lemma}{original lemmas}
\Crefname{originallemma}{Original lemma}{Original lemmas}
\makeatletter
\newlabel{categories-lemma-characterize-left-exact}{{0035}{source}{}{sourcepointer-7}{}}
\newlabel{categories-lemma-characterize-left-exact@cref}{{[originallemma][0][]0035}{[1][0][]source}{}{}{}}
\makeatother
\crefname{originallemma}{original lemma}{original lemmas}
\Crefname{originallemma}{Original lemma}{Original lemmas}
\makeatletter
\newlabel{lemma-compute-j-shriek-restrict}{{03EE}{source}{}{sourcepointer-8}{}}
\newlabel{lemma-compute-j-shriek-restrict@cref}{{[originallemma][0][]03EE}{[1][0][]source}{}{}{}}
\makeatother
\title{Localizing a morphism of sites}
\author{The Stacks Project Authors}
\date{stacks-project:sites.tex:local-snapshot-20261012; GNU FDL 1.2 or later}
\begin{document}
\maketitle
\noindent\textit{Editorial note.} All delivered mathematical text is the author's own, copied from the pinned Stacks Project chapter file sites-cohomology.tex; no statement or proof step was rewritten, reordered or removed.\par\medskip
\noindent\textit{Editorial note.} The proof is the author's complete argument; every step the author writes out is inside the counted span and no closing sentence is a placeholder.\par\medskip
\noindent\textit{Editorial note.} Every labelled reference occurring inside the counted spans is carried in assembly.external\_references as a non-load-bearing navigation pointer (same-chapter ones at their definition site, cross-chapter ones at the invocation inside the counted span), so no load-bearing mathematics is replaced by a stub. No bibliography citation occurs inside the counted spans.\par\medskip
\noindent\textit{Editorial note.} The checked chapter list for same-name collisions is the set of stacks chapters staged in this repository: chow, cohomology, divisors, etale-cohomology, flat, more-algebra, more-morphisms, morphisms, spaces-more-morphisms, varieties together with this card's chapter sites-cohomology.\par\medskip
\noindent\textit{Editorial note.} The preamble is the chapter's own retained minimal preamble (sources/stacks-sites-cohomology/retained-preamble.tex), a verbatim excerpt of the native preamble: the Xy-pic setup and the native theorem declarations with the mathematical macros.\par\medskip
\noindent\textit{Allowed standard prerequisites (editorial).}\begin{itemize}
\item The equivalence between sheaves on the localized site and sheaves on a slice. Same chapter, Lemma lemma-essential-image-j-shriek. Identifies the categories of sheaves on the localized sites with the slices, which is what makes the description of the pullback meaningful.
\item The computation of the lower shriek of a localization. Same chapter, Lemma lemma-technical-up. Supplies the identifications used to compare the composites of lower shriek and pullback.
\item The pullback of a representable sheaf. Same chapter, Lemma lemma-pullback-representable-sheaf. Gives the identification of the pullback of the sheaf represented by an object.
\end{itemize}\medskip
\ifdefined\hypertarget\hypertarget{sourcepointer-1}{}\fi
\noindent\textit{Source reference.} Original lemma 00Y1: sources/stacks-sites/pinned-original-source.tex, source lines 6438--6438. Navigation pointer invoked at proof line 6438; the counted span contains the pointer only and the cited result is not reproduced inside it, so the pointer is registered as a non-load-bearing navigation stub in the established form for Stacks material. No bibliography citation is involved.\par\medskip
\ifdefined\hypertarget\hypertarget{sourcepointer-2}{}\fi
\noindent\textit{Source reference.} Original section 00VC: sources/stacks-sites/pinned-original-source.tex, source lines 6514--6514. Navigation pointer invoked at proof line 6514; the counted span contains the pointer only and the cited result is not reproduced inside it, so the pointer is registered as a non-load-bearing navigation stub in the established form for Stacks material. No bibliography citation is involved.\par\medskip
\ifdefined\hypertarget\hypertarget{sourcepointer-3}{}\fi
\noindent\textit{Source reference.} Original section 00WU: sources/stacks-sites/pinned-original-source.tex, source lines 6515--6515. Navigation pointer invoked at proof line 6515; the counted span contains the pointer only and the cited result is not reproduced inside it, so the pointer is registered as a non-load-bearing navigation stub in the established form for Stacks material. No bibliography citation is involved.\par\medskip
\ifdefined\hypertarget\hypertarget{sourcepointer-4}{}\fi
\noindent\textit{Source reference.} Original lemma 00WY: sources/stacks-sites/pinned-original-source.tex, source lines 6543--6543. Navigation pointer invoked at proof line 6543; the counted span contains the pointer only and the cited result is not reproduced inside it, so the pointer is registered as a non-load-bearing navigation stub in the established form for Stacks material. No bibliography citation is involved.\par\medskip
\ifdefined\hypertarget\hypertarget{sourcepointer-5}{}\fi
\noindent\textit{Source reference.} Original lemma 04D3: sources/stacks-sites/pinned-original-source.tex, source lines 6560--6560. Navigation pointer invoked at proof line 6560; the counted span contains the pointer only and the cited result is not reproduced inside it, so the pointer is registered as a non-load-bearing navigation stub in the established form for Stacks material. No bibliography citation is involved.\par\medskip
\ifdefined\hypertarget\hypertarget{sourcepointer-6}{}\fi
\noindent\textit{Source reference.} Original lemma 0038: sources/stacks-sites/pinned-original-source.tex, source lines 6569--6569. Navigation pointer invoked at proof line 6569; the counted span contains the pointer only and the cited result is not reproduced inside it, so the pointer is registered as a non-load-bearing navigation stub in the established form for Stacks material. No bibliography citation is involved.\par\medskip
\ifdefined\hypertarget\hypertarget{sourcepointer-7}{}\fi
\noindent\textit{Source reference.} Original lemma 0035: sources/stacks-sites/pinned-original-source.tex, source lines 6573--6573. Navigation pointer invoked at proof line 6573; the counted span contains the pointer only and the cited result is not reproduced inside it, so the pointer is registered as a non-load-bearing navigation stub in the established form for Stacks material. No bibliography citation is involved.\par\medskip
\ifdefined\hypertarget\hypertarget{sourcepointer-8}{}\fi
\noindent\textit{Source reference.} Original lemma 03EE: sources/stacks-sites/pinned-original-source.tex, source lines 6581--6581. Navigation pointer invoked at proof line 6581; the counted span contains the pointer only and the cited result is not reproduced inside it, so the pointer is registered as a non-load-bearing navigation stub in the established form for Stacks material. No bibliography citation is involved.\par\medskip
\begin{lemma}
\label{lemma-localize-morphism}
Let $f : \mathcal{C} \to \mathcal{D}$ be a morphism of sites
corresponding to the continuous functor $u : \mathcal{D} \to \mathcal{C}$.
Let $V \in \Ob(\mathcal{D})$ and set $U = u(V)$.
Then the functor $u' : \mathcal{D}/V \to \mathcal{C}/U$,
$V'/V \mapsto u(V')/U$ determines a morphism of sites
$f' : \mathcal{C}/U \to \mathcal{D}/V$.
The morphism $f'$ fits into a commutative diagram of topoi
$$
\xymatrix{
\Sh(\mathcal{C}/U) \ar[r]_{j_U} \ar[d]_{f'} &
\Sh(\mathcal{C}) \ar[d]^f \\
\Sh(\mathcal{D}/V) \ar[r]^{j_V} &
\Sh(\mathcal{D}).
}
$$
Using the identifications
$\Sh(\mathcal{C}/U) = \Sh(\mathcal{C})/h_U^\#$ and
$\Sh(\mathcal{D}/V) = \Sh(\mathcal{D})/h_V^\#$ of
Lemma \ref{lemma-essential-image-j-shriek}
the functor $(f')^{-1}$ is described by the rule
$$
(f')^{-1}(\mathcal{H} \xrightarrow{\varphi} h_V^\#)
=
(f^{-1}\mathcal{H} \xrightarrow{f^{-1}\varphi} h_U^\#).
$$
Finally, we have $f'_*j_U^{-1} = j_V^{-1}f_*$.
\end{lemma}

\begin{proof}
It is clear that $u'$ is continuous, and hence we get functors
$f'_* = (u')^s = (u')^p$ (see
Sections \ref{section-functoriality-PSh}
and \ref{section-continuous-functors})
and an adjoint $(f')^{-1} = (u')_s = ((u')_p\ )^\#$. The assertion
$f'_*j_U^{-1} = j_V^{-1}f_*$ follows as
$$
(j_V^{-1}f_*\mathcal{F})(V'/V)
= f_*\mathcal{F}(V') = \mathcal{F}(u(V'))
= (j_U^{-1}\mathcal{F})(u(V')/U)
= (f'_*j_U^{-1}\mathcal{F})(V'/V)
$$
which holds even for presheaves. What isn't clear a priori is
that $(f')^{-1}$ is exact, that the diagram commutes, and that
the description of $(f')^{-1}$ holds.

\medskip\noindent
Let $\mathcal{H}$ be a sheaf on $\mathcal{D}/V$.
Let us compute $j_{U!}(f')^{-1}\mathcal{H}$. We have
\begin{eqnarray*}
j_{U!}(f')^{-1}\mathcal{H}
& =
((j_U)_p(u'_p\mathcal{H})^\#)^\# \\
& =
((j_U)_pu'_p\mathcal{H})^\# \\
& =
(u_p(j_V)_p\mathcal{H})^\# \\
& =
f^{-1}j_{V!}\mathcal{H}
\end{eqnarray*}
The first equality by unwinding the definitions.
The second equality by Lemma \ref{lemma-technical-up}.
The third equality because $u \circ j_V = j_U \circ u'$.
The fourth equality by Lemma \ref{lemma-technical-up} again.
All of the equalities above are isomorphisms of functors, and
hence we may interpret this as saying that the following
diagram of categories and functors is commutative
$$
\xymatrix{
\Sh(\mathcal{C}/U) \ar[r]_{j_{U!}} &
\Sh(\mathcal{C})/h_U^\# \ar[r] &
\Sh(\mathcal{C}) \\
\Sh(\mathcal{D}/V) \ar[r]^{j_{V!}} \ar[u]^{(f')^{-1}} &
\Sh(\mathcal{D})/h_V^\# \ar[r] \ar[u]^{f^{-1}} &
\Sh(\mathcal{D}) \ar[u]^{f^{-1}}
}
$$
The middle arrow makes sense as $f^{-1}h_V^\# = (h_{u(V)})^\# = h_U^\#$, see
Lemma \ref{lemma-pullback-representable-sheaf}.
In particular this proves the description of $(f')^{-1}$ given
in the statement of the lemma.
Since by
Lemma \ref{lemma-essential-image-j-shriek}
the left horizontal arrows are equivalences
and since $f^{-1}$ is exact by assumption we conclude that
$(f')^{-1} = u'_s$ is exact. Namely, because it is a left adjoint
it is already right exact
(Categories, Lemma \ref{categories-lemma-adjoint-exact}).
Hence we only need to show that
it transforms a final object into a final object and commutes
with fibre products
(Categories, Lemma \ref{categories-lemma-characterize-left-exact}).
Both are clear for the induced functor
$f^{-1} : \Sh(\mathcal{D})/h_V^\# \to \Sh(\mathcal{C})/h_U^\#$.
This proves that $f'$ is a morphism of sites.

\medskip\noindent
We still have to verify that $(f')^{-1}j_V^{-1} = j_U^{-1}f^{-1}$.
To see this use the formula above and the description
in Lemma \ref{lemma-compute-j-shriek-restrict}. Namely,
combined these give, for any sheaf $\mathcal{G}$ on $\mathcal{D}$, that
$$
j_{U!}(f')^{-1}j_V^{-1}\mathcal{G}
=
f^{-1}j_{V!}j_V^{-1}\mathcal{G}
=
f^{-1}(\mathcal{G} \times h_V^\#)
=
f^{-1}\mathcal{G} \times h_U^\#
=
j_{U!}j_U^{-1}f^{-1}\mathcal{G}.
$$
Since the functor $j_{U!}$ induces an equivalence
$\Sh(\mathcal{C}/U) \to \Sh(\mathcal{C})/h_U^\#$
we conclude.
\end{proof}

\end{document}
